\chapter{Le formatage et le partitionnement}
	\section{Définition du formatage}
	Le formatage consiste à préparer un support de stockage en y inscrivant un système de fichiers afin qu'il soit reconnu par le système d'exploitation.

Parmi les systèmes de fichiers les plus courants figurent :
\begin{itemize}[label=$\bullet$]
    \item FAT
    \item NTFS
    \item HFS
    \item ext2
    \item ext3
    \item UFS
\end{itemize}

	\section{Le partitionnement}
Les disques de grande capacité peuvent être divisés en plusieurs partitions logiques. Chaque partition peut recevoir un système de fichiers spécifique.

Cette opération est appelée partitionnement.

	\section{Les différents niveaux de formatage}
	\subsection{Le formatage de bas niveau}
	Le formatage de bas niveau prépare la surface physique du disque afin qu'elle soit conforme aux attentes du contrôleur matériel.

Cette opération est généralement réalisée par le fabricant.
	\subsection{Le formatage de bas niveau}
Le formatage de haut niveau consiste à inscrire les informations logicielles correspondant au système de fichiers choisi.

Cette opération est effectuée par l'utilisateur ou le système d'exploitation.